\subsection{二重数列と二方向の階差}\label{節：二重数列と二方向の階差}
  一つの番号に沿って項を並べる代わりに,行と列の二つの番号で値を並べてみよう.
  ここでは$n,m\ge0$を整数とし,$a_{n,m}$を数表の位置$(n,m)$の値とする.
  このように二つの添字をもつものを\textbf{二重数列}という.
  新しい万能の解法ではなく,階差と帰納法を二方向で使う見方を学ぶ.

  \subsubsection{距離を二乗し,二方向へ差を取る}
  原点から点$(n,m)$までの距離は
  \[ a_{n,m}=\sqrt{n^2+m^2} \]
  である.
  根号を取り除くため$b_{n,m}=a_{n,m}^2$とすると,
  一方向の差は$b_{n+1,m}-b_{n,m}=2n+1$となり,$m$によらない.
  従って,さらにもう一方向へ差を取れば零になる.

  この計算を広げ,境界の値
  \[
    b_{n,0}=n^2,\qquad b_{0,m}=m^2
  \]
  と,一定の実数$c$について
  \BookMathLayout{\[
    b_{n+1,m+1}-b_{n+1,m}-b_{n,m+1}+b_{n,m}=c
  \]}{\[
    \begin{aligned}
      b_{n+1,m+1}-b_{n+1,m}\\
      {}-b_{n,m+1}+b_{n,m}=c
    \end{aligned}
  \]}
  が与えられているとする.内側の値はこの式と境界から順に決まる.
  右辺を$b_{n+1,m+1}$について解けば,ほかの三つの位置は添字の和が小さい.
  $n+m$についての帰納法により一意性を確かめられる.

  横方向の差$d_{n,m}=b_{n+1,m}-b_{n,m}$を追うと,
  \[
    d_{n,m+1}-d_{n,m}=c,\qquad d_{n,0}=2n+1
  \]
  なので$d_{n,m}=2n+1+cm$である.これを横に足し戻せば,
  \[
    b_{n,m}=n^2+m^2+cnm
  \]
  となる.\textbf{二方向に一つずつ階差を取り,逆順に足し戻した}のである.
  $c$が一定という条件が,この短い計算を可能にしている.

  \subsubsection{図形と数え上げに戻る}
  二つの単位ベクトル$\boldsymbol u,\boldsymbol v$のなす角を$\theta$とすると,
  $c=2\cos\theta$の場合には
  \[
    b_{n,m}=|n\boldsymbol u+m\boldsymbol v|^2
  \]
  と読める.この図形的な解釈には$-2\le c\le2$が必要である.
  単に数表を作る式としてなら$c$は任意の実数でよいが,
  その値を距離の二乗と読むためには非負性も確認する.

  一方,格子点$(0,0)$から$(n,m)$へ,右と上へ一歩ずつ進む道の数を$A_{n,m}$とする.
  最後の一歩を分ければ,
  \[
    A_{n,m}=A_{n-1,m}+A_{n,m-1}\quad(n,m\ge1)
  \]
  であり,境界では$A_{n,0}=A_{0,m}=1$である.
  \BookFigLatticeRecurrence
  下の数表は,上と左の値を足して作ったものである.
  \[
    \begin{array}{c|rrrr}
      m\backslash n&0&1&2&3\\\hline
      0&1&1&1&1\\
      1&1&2&3&4\\
      2&1&3&6&10\\
      3&1&4&10&20
    \end{array}
  \]
  全$n+m$歩から上へ進む$m$歩の位置を選ぶので,
  \[
    A_{n,m}=\binom{n+m}{m}
  \]
  である.最後の一歩で分ける説明は,二項係数の加法公式の説明でもある.
  この公式と境界を使った帰納法による確認は,
  \bookproblemref{practice-connections}{3}{接続演習問3}で行おう.

  \begin{bookpoint}
    \textbf{添字を増やしても,確認することは同じ}\par
    どこが初期条件に当たる境界か,どの値から次が決まるかを先に確かめる.
    二方向の関係を別々に与える場合は,同じ位置を二通りに計算して矛盾しないことも必要である.
  \end{bookpoint}
  格子経路と二項係数の資料は参考文献のNIST DLMF §26.3へ.
